% ============================================================================
% HEC Montréal MBA -- ECON50803: Environnement macroéconomique
% Shared Beamer Preamble
% ============================================================================
% Loads jfb-teaching-slides.sty from latex-common (sibling repo) for the full
% teaching Beamer aesthetic: Fira Sans fonts, HEC palette, titleframe /
% sectionframe, box envs (keyinsight, bizimplication, warning, defbox,
% exercise), newspaper headlines, backup slides, TikZ styles for AD-AS and
% supply/demand diagrams.
%
% Course-specific configuration, project paths, and overrides live below.
% ============================================================================

\usepackage{jfb-teaching-slides}

% --- Course metadata (drives the footer) ------------------------------------
\jfbcoursename{ECON 50803}
\jfbcoursetitle{Environnement macroéconomique}
\jfbinstitution{HEC Montréal}

% --- Project paths ----------------------------------------------------------
\graphicspath{{../Figures/}}
\newcommand*\includetables[1]{\input{../Tables/#1.tex}}

% --- Shared closing roadmap -----------------------------------------------
% The argument is the session just completed. On the white slide background,
% HECnavy!40 gives completed themes the appearance of 40% opacity.
\newcommand{\coursethemestyle}[3]{%
   \ifnum#1>#2\relax
      \textcolor{HECnavy}{#3}%
   \else
      \textcolor{HECnavy!40}{#3}%
   \fi
}
\newcommand{\courseroadmap}[1]{%
   \begin{frame}{Les thèmes du cours}
      \renewcommand{\arraystretch}{1.4}
      \centering
      \small
      \begin{tabularx}{\textwidth}{>{\bfseries\raggedright\arraybackslash}X >{\raggedright\arraybackslash}p{0.40\textwidth}}
         \toprule
         \textbf{Question centrale} & \textbf{Séance} \\
         \midrule
         \coursethemestyle{1}{#1}{Comment mesurer la production, les prix et le niveau de vie\,?} & \coursethemestyle{1}{#1}{S1 : Mesurer l'économie} \\
         \coursethemestyle{2}{#1}{Pourquoi certains pays sont-ils plus riches que d'autres\,?} & \coursethemestyle{2}{#1}{S2 : Croissance de long terme} \\
         \coursethemestyle{3}{#1}{Qu'est-ce qui détermine l'emploi, les salaires et l'investissement\,?} & \coursethemestyle{3}{#1}{S3 : Marchés du travail et du capital} \\
         \coursethemestyle{4}{#1}{Pourquoi l'économie connaît-elle des hauts et des bas\,?} & \coursethemestyle{4}{#1}{S4 : Cycles économiques} \\
         \coursethemestyle{5}{#1}{Pourquoi et comment combattre les récessions\,?} & \coursethemestyle{5}{#1}{S5 : Politique monétaire et budgétaire} \\
         \bottomrule
      \end{tabularx}
   \end{frame}
}

% --- TikZ style overrides (preserve pre-migration visual choices) ----------
% The shared jfb-teaching-slides.sty defines stronger defaults (e.g.
% supply/demand in HECnavy, eqpoint default HECgreen). Keep the original
% ECON50803 overrides here so existing slides render unchanged.
\tikzset{
   supply/.style     = {very thick},
   original/.style   = {HECgray, thick, dashed},
   refline/.style    = {thin, dashed, HECgray},
   shiftarrow/.style = {very thick, -Stealth},
   eqpoint/.default  = HECnavy,
   eqfaded/.style    = {fill=HECgray, draw=HECgray, opacity=0.4},
   potline/.style    = {HECgreen, thick, dashed},
   econbrace/.style  = {decorate, decoration={brace, amplitude=6pt, mirror}, thick},
   pointlabel/.style = {font=\small},
   annotlabel/.style = {font=\small},
}

% --- Exam-relevance tag -----------------------------------------------------
% Set \showexamtagstrue and rebuild all sessions before the exam.
\newif\ifshowexamtags
\showexamtagsfalse
\usetikzlibrary{shapes.geometric}
\newcommand{\examtag}{%
   \ifshowexamtags%
      \begin{tikzpicture}[remember picture, overlay]
         \node[star, star points=5, star point ratio=2.25,
               fill=HECcoral, inner sep=1.6pt, anchor=north east]
            at ([xshift=-0.5cm, yshift=-0.65cm]current page.north east) {};
      \end{tikzpicture}%
   \fi%
}
